\chapter{Toothache (Dental Pain)}

\section*{Definition}
Pain originating from a tooth or its supporting structures, ranging from mild sensitivity to severe throbbing pain, caused by dental caries, infection, trauma, or dental procedures.

\section*{Pathophysiology}
Dental caries (decay) progresses through enamel into dentin, exposing sensitive nerve endings. Pulpitis (inflammation of dental pulp) results from deep decay, trauma, or thermal injury. Periapical abscess forms when pulp infection extends through the apical foramen into surrounding bone. Periodontal abscess arises from gum infection in periodontal pockets.

\section*{Etiology}
\begin{itemize}
  \item Dental caries (tooth decay) — most common cause
  \item Pulpitis (reversible or irreversible)
  \item Periapical abscess (dental abscess)
  \item Periodontal abscess (gum infection)
  \item Cracked tooth syndrome
  \item Dental trauma (fractured or avulsed tooth)
  \item Impacted wisdom teeth (pericoronitis)
  \item Temporomandibular joint (TMJ) disorders (referred pain)
  \item Sinusitis (referred pain to upper teeth)
  \item Dental procedures (recent filling, root canal, extraction)
\end{itemize}

\section*{Indications for Evaluation}
Seek care if:
\begin{itemize}
  \item Severe or worsening symptoms
  \item Severe tooth pain with facial swelling, fever, difficulty swallowing or breathing (possible deep space neck infection), or persistent pain beyond 2 days
  \item Accompanied by other concerning signs
\end{itemize}

\section*{Related Symptoms}
\begin{itemize}
  \item Facial swelling
  \item Fever
  \item Headache
  \item Gum swelling
  \item Earache (referred)
\end{itemize}
\textbf{Note:} Always consider serious underlying conditions when evaluating persistent toothache (dental pain).

\section*{Self-Care / Home Management}
\begin{itemize}
  \item Rest and avoid activities that may aggravate the condition.
  \item Maintain adequate hydration and a balanced diet.
  \item Over-the-counter remedies may provide symptomatic relief (consult a pharmacist or doctor).
  \item Monitor symptoms and seek professional advice if they persist or worsen.
  \item Avoid self-medicating with prescription medications without medical guidance.
\end{itemize}

\section*{Diagnostic Approach}
\begin{itemize}
  \item A thorough history and physical examination are the first steps in evaluation.
  \item Laboratory tests (blood work, urinalysis) may be ordered based on clinical suspicion.
  \item Imaging studies (X-ray, ultrasound, CT, MRI) may be indicated when structural pathology is suspected.
  \item Specialist referral is arranged when the diagnosis remains uncertain or requires specific expertise.
\end{itemize}

\section*{Treatment / Management Overview}
\begin{itemize}
  \item Treatment targets the underlying cause identified through diagnostic evaluation.
  \item Supportive care and symptom management are provided while awaiting definitive therapy.
  \item Medications, lifestyle modifications, and physical therapies may be prescribed as appropriate.
  \item Follow-up ensures treatment effectiveness and allows timely adjustments to the care plan.
\end{itemize}

\clearpage